\thispagestyle{plain}
\ctsection{Préface de la nouvelle édition}

\begin{otherlanguage}{french}
\begin{multicols}{2}
\textit{En réimprimant ce livre, fameux en Allemagne au siècle dernier, notre unique but a été de rendre aux vrais amis des études latines un ouvrage destiné à développer la connaissance du latin. Nous avons voulu offrir aux professeurs et aux élèves, surtout à ces derniers, un manuel vraiment pratique, pour approfondir les beautés intimes de la langue qui “ seule entre toutes les langues mortes, est véritablement ressuscitée, et qui UNE FOIS RESSUSCITÉE, NE MEURT PLUS!”}\footnote{de Maistre, \textit{Du Pape}.}

\textit{Avant d’indiquer le plan du livre et la manière de s’en servir utilement, qu’il nous soit permis de rappeler en quelques mots les principaux actes de la vie de l’auteur.}

\textit{ François Wagner naquit à Wangen, en Souabe, le 14 août 1675. Il fut admis dans la Compagnie de Jésus au noviciat de Krems à l’âge de 15 ans. Après avoir terminé ses deux années d’épreuve, il enseigna pendant plus de 25 ans la rhétorique, dans les colléges de Krems, de Presbourg et de Tyrnau. Étant recteur du Séminaire de Vienne, il dirigea une congrégation et s’occupa de la composition de plusieurs ouvrages estimés. Il mourut à Vienne, le 8 février 1738, à l’âge de 62 ans.}

\textit{Tels sont les détails que nous fournit sur cet homme illustre la Bibliothèque des écrivains de la Compagnie de Jésus. Le P. Jean Nép. Stœger, dans ses Scriptores Provinciæ Austriacæ Societatis Jesu, nous a transmis la liste complète des nombreux ouvrages écrits par le Père Wagner, qui tour à tour, s’occupant de rhétorique, de grammaire, de géographie, d’histoire, nous a laissé dans ces diverses branches des connaissances humaines, des traces remarquables de son esprit actif et pénétrant.}

\textit{Notre but, dans cet avertissement, est de parler seulement de son Dictionnaire de Phraséologie. Mais laissons l’auteur lui-même nous raconter comment naquit ce livre, nous indiquer ce qu’il a voulu en le composant, et de quelle manière il est parvenu à mener à terme l’ouvrage que nous possédons.}\footnote{Voici, d’après les Auteurs de la \textit{Bibliothèque des Écrivains de la Compagnie de Jésus}, Tome 3ᵉ, pages 1467, 1468 et 1469, (Louvain, MDCCCLXXVI) les principales éditions de l’ouvrage que nous offrons aujourd’hui au public:\par Universæ Phraseologiæ corpus congestum. Augustæ Vindelicorum, 1718.—Ratisbonæ, 1740, 8ᵒ.\par Universæ Phraseologiæ corpus congestum a P. Francisco Wagner, Societatis Jesu Sacerdote, secundis curis a quopiam ejusdem Societatis Sallustiana, Cæsareana, Liviana, Corneliana, etc., Phraseologiis locupletatum. Editio novissima, auctior et emendatior. Cum Privilegio Sacræ Cæsareæ Majestatis et Permissu Superiorum. Ratisbonæ, Sumptibus Emerici Felicis Baderi, A. P. C. N. MDCCXLV, 8ᵒ, \textit{pp.} 816 \textit{à 2 col.}, sll. \textit{L’épître dédicatoire, signée par l’auteur, est en date du 20 février 1718. La concession à l’imprimeur est du 4 janvier 1738. A la fin on a joint}: Latinitatis ars brevicula, seu syntaxis ornata, de tribus latinæ linguæ virtutibus, puritate, elegantia et copia, \textit{pp.} 82. Index diffusior vocum barbararum vel minus elegantium, \textit{p.} 1 \textit{à} 48. Index Germanico-Latinus \textit{p.} 49 \textit{à} 96.—\textit{Même titre.} Ratisbonæ, Sumptibus Emerici Felicis Baderi, A. P. C. N. MDCCLVIII, 8ᵒ, \textit{pp.} 808, 80 \textit{et} 85—\textit{Ibid.} id. 1751, 8ᵒ.\par Universæ Phraseologiæ Latinæ corpus congestum a P. Francisco Wagner, Societatis Jesu Sacerdote, secundis curis a quopiam ejusdem Societatis Sallustiana, Cæsareana, Liviana, Corneliana, etc., Phraseologiis ac denique Indice Verborum quæ in foro militari, civili sacroque obtinent, locupletatum. Editio Novissima, Auctior et Emendatior. Ratisbonæ et Viennæ, Sumptibus Emerici Felicis Baderi, 1760, 8ᵒ, \textit{pp.} 808 \textit{à} 2 \textit{col.} sll. \textit{A la fin on a joint}: Latinitatis ars brevicula, seu Syntaxis ornata, de tribus Latinæ linguæ virtutibus, puritate, elegantia et copia, \textit{pp.} 80. Index diffusior vocum barbararum, aut minus elegantium, \textit{p.} 1 \textit{à} 48. Index vocum quæ in foro militari, civili, sacroque obtinent. Ex Iano, Noltenio, d’Aquino, Seyboldo, Strada, Maffejo, Galutio, \textit{etc.}, collectus, \textit{p.} 49 \textit{à} 74. Index Germanico-Latinus, \textit{p.} 75 \textit{à} 111, \textit{fin du volume}.\par P. Francisci Wagner Universa Phraseologia Latina. Ab eodem secundis curis Sallustiana, Cæsareana, Liviana, Corneliana, etc. Phraseologiis ac denique Indice Verborum quæ in foro militari, civili sacroque obtinent, locupletata et ad usum juventutis litterarum studiosæ accommodata. Editio novissima, auctior et emendatior. Augustæ Vindelicorum, Sumptibus Matthæi Rieger p. m. filiorum. MDCCXCII, 8ᵒ, \textit{pp.} XVI \textit{à} 896. \textit{A la fin}: Latinitatis ars brevicula, \textit{etc.}, \textit{pp.} 204.\par Stœger dit: “Universæ.... Corpus auctum ab alio Sacerdote, S. J. Phraseologia Sallust. Cæsar. Liv. Cornel. etc., Augustæ Vindelicorum, 1801, 8ᵒ, et cum indice Viennæ, Geistinger, 1824, 8ᵒ.”\par Fr. Wagner Universa Phraseologia Latina, edit. novissima, Schmidt, Viennæ et Tergesti, 1824, 8ᵒ.\par Wagneri Universa Phraseologia Latina, aucta et emendata a Span. Viennæ, 1825.\par Adjecta Phraseologia Germanico-Latina a P. Goldhagen, S. J. cum Indice Verborum in foro sacro, civili et militari obtinentium. Moguntiæ, 1751, 1766, 8ᵒ.\par Wagner’s deutsch-lateinische Phraseologie, bearbeitet von Seibt, Prag. 1846.\par Universa Phraseologia Latina adjecta lingua hungarico slavonica. Tyrnaviæ, 1775, 8ᵒ.\par Franc. Wagneri Universæ Phraseologiæ corpus idiomate slavico locupletatum. Tyrnaviæ, 1750, 8ᵒ \textit{Par le P. B.} Zambal, S. J.\par Compendiaria methodus addiscendi tres præcipuas latinæ linguæ virtutes, Puritatem, Elegantiam et Copiam, in usum juventutis emendatæ latinitatis studiosæ proposita. Authore R. P. Francisco Wagner Societatis Jesu. Cum facultate superiorum. Dilingæ, Sumptibus Wilhelmi Hausen, anno 1735, 8ᵒ, \textit{pp.} 96.\par Supellex latinitatis ex Phraseologia Francisci Wagner, S. J. ad usum scholarum ejusdem Societatis collecta. Vilnæ, typ. Academ. S. J., 1751, 8ᵒ. (Editor est P. Fr. Bohomolec S. J., teste Ianocki.)—Idem, cum vocibus polonis. Calissii, typ. S. J., 1752, 12ᵒ.—Editio 2ᵃ Leopoli, typ. S. J., 1756, 8ᵒ, \textit{pp.} 42, 377 \textit{et} 54.—Vilnæ, typ. Academ. S. J., 1757, 12ᵒ.—Calissii, typ. S. J., 1772, 12ᵒ.\par Supellex latinitatis ex Phraseologia, P. Fr. Wagner S. J., ad usum scholarum ejusdem Societatis collecta. Editio 2ᵃ correctior. Leopoli, typ. S. J., 1756, 8ᵒ, \textit{pp.} 23 \textit{à} 77 \textit{et} 102.\par Outre ces 24 éditions, nous en possédons deux autres. La première appartient à la bibliothèque de St Acheul: Universæ, etc... Corpus, demum linguis Hungarica, Germanica et Slavica locupletatum. Budæ, typis typographiæ Regiæ Univers. Hungaricæ, 1822, 8ᵒ, 2 vol. \textit{pp.} 1524 à 83. L’autre, appartenant à la bibliothèque de l’École libre de la Providence, est un peu plus ancienne. C’est une réimpression de l’édition de 1792, faite chez Mathias Rieger, \textit{pp.} XVI \textit{à} 896; supp. 204. Augustæ Vindelicorum, MDCCC.}

\textit{“Étant professeur des jeunes religieux qui recommençaient leurs humanités, le but de tous mes soins était d’accomplir le sage précepte du P. Jouvency, sans contredit le premier parmi nous dans la connaissance parfaite de la langue latine,” Patris Juvencii, harum litterarum hodie inter nos facile Principis. “Le premier travail du maître, dit celui-ci, sera de montrer à ses élèves le style le plus pur et de les aider à l’acquérir. Je recherchais donc ce que je pouvais offrir de mieux à ces jeunes gens, pour les conduire aussi promptement que sûrement à l’acquisition des trois grandes qualités de la langue latine, la propriété, l’élégance, l’abondance. Pour la propriété, outre les Thesaurus, je faisais étudier tous les jours à mes élèves pendant un quart d’heure le Flos Latinitatis du P. Pomey, et au bout de deux mois, je remarquais des fruits merveilleux, jam intra bimestre Latinum quid velut in herba efflorescere, non sine jucundo animi sensu animadverti. Pour l’élégance, Tursellini, Schorus, Dolet, Hadrianus étaient les maîtres que je leur faisais étudier et copier. Enfin, pour l’abondance, Manusius, Bucher et Schönsleder. Mais, bientôt, je m’aperçus que cette foule de livres inondaient leurs pupitres, qu’ils perdaient un temps précieux à rechercher, ici les adverbes, là les épithètes ou les synonymes, ailleurs les développements, et c’est alors que je résolus de réunir avec ordre et méthode, dans un seul ouvrage, le fruit des veilles de tous ces hommes illustres.”}

\textit{C’est ainsi que naquit le Universæ Phraseologiæ corpus congestum. Il fut d’abord exclusivement composé de phrases et de tournures tirées de Cicéron, puis on y ajouta un petit nombre d’expressions appartenant aux grands auteurs du siècle d’or de la langue latine, Salluste, César, Tite-Live, Cornelius et quelques autres. Il faut remarquer que l’auteur ne cite pas toujours textuellement l’écrivain auquel il a emprunté sa citation, et c’est ce qui a empêché, dans cette édition, d’indiquer les sources, quoique nous puissions affirmer que nous avons vérifié presque toutes les citations au moyen de l’Apparatus de Nizzolius, de Forcellini et de Freund. Souvent il se contente de prendre la tournure et de l’appliquer au passage qu’il veut développer. Cependant, dans ce qui est rangé sous la rubrique Usus : les phrases qui servent à indiquer la Syntaxe particulière de chaque mot, sont empruntées textuellement à Cicéron et aux meilleurs auteurs du grand siècle. Souvent nous les avons rétablies et traduites, pour la plus grande utilité des élèves. “Leur lecture assidue, ajoute le P. Wagner, non-seulement donnera aux élèves les différentes significations d’un mot, mais encore leur prouvera que négliger d’en suivre la construction, c’est s’exposer à employer une tournure peu latine, “ si quampiam atque ab his alienam construendi formulam adhibuerint Discipuli, fere in Idiotismum se impactos admoneantur.”}

\textit{Mais il ne suffit pas d’avoir un bon livre, il faut savoir s’en servir; or, à première vue, il semble que l’usage de ce Dictionnaire ne soit pas accessible à tout le monde. Il est hors de doute, et l’expérience le prouve, que si l’on parvient à en saisir l’économie, si l’on en a, pour ainsi dire, la clef, on parviendra en peu de temps à faire de rapides progrès dans la connaissance de la langue latine. L’auteur en avait constaté l’efficacité, quand il nous dit: “ Noveram adolescentes qui orationem struerent tam elegantem, ut in volvendo Cicerone, multos annos posuisse videri potuissent hujus artificii ignari. Inest et illud commodi quod ea, quæ aliena fuerant, multiplici usu sua faciant; quodque dicitur, in succum et sanguinem suum vertant. ”}

\textit{Pour cela, il ne faut pas puiser indistinctement dans ce livre, coudre ensemble, sans suite et sans liens, des membres de phrase, qui, pris individuellement, sont très latins, mais qui doivent être rapprochés avec goût et discrétion. L’enfant intelligent doit faire sans dictionnaire son discours, sa narration ou son thème; il doit, après avoir médité profondément son sujet, après l’avoir ordonné dans toutes ses parties, l’écrire dans un latin plus ou moins pur, plus ou moins élégant, selon sa force et sa capacité. Ce travail, le plus important sans contredit, étant terminé, il devra, avec l’aide du dictionnaire de Phraséologie, débarrasser son devoir, que nous pouvons supposer correct, de toutes les tournures françaises, des gallicismes, des mots peu latins, et de cent autres imperfections, auxquelles il trouvera un remède assuré dans ce volume, s’il sait s’en servir avec intelligence et discernement.}

\textit{Montrons, par un exemple facile, la vraie manière d’user avec fruit de ce livre. C’est une lettre d’un fils à son père pour lui demander de l’argent, dans le but d’acheter des livres. Dans la première colonne, se trouve cette lettre dans le style d’un élève ordinaire de troisième; style simple et exempt de fautes contre les règles de la grammaire. Dans la seconde, la même lettre a été modifiée, rendue plus élégante au moyen du dictionnaire de Phraséologie. Du reste, les mots modifiés dans la première colonne sont soulignés, ce sont ceux-là qui doivent être l’objet du travail personnel de l’élève, quand il a jeté sur le papier les idées et qu’il veut donner à son style ce poli, cette perfection que demande la langue latine.}

\end{multicols}
\noindent\begin{minipage}[t]{0.48\textwidth}\small
\textit{Scio} equidem, te alioquin mei causa multos \textit{sumptus facere}; ausus tamen sum, novam insuper gratiam \textit{petere}. \textit{Suadet} nobis magister, ut libros quos in schola \textit{explicat, coemeremus, spondetque}, si eos domi \textit{legerimus}, multum profecturos. Uti hactenus semper \textit{conatus sum}, primo \textit{loco} aut primis proximus esse, ita nihil \textit{tristius} foret, quam si, ob hunc librorum defectum, priorem laudem \textit{amitterem} et a discipulis \textit{contemnerer}. \textit{Spero} autem in causa tam æqua, te, ut fuisti, erga me \textit{liberalem} fore.


\end{minipage}\hfill
\begin{minipage}[t]{0.48\textwidth}\small
Cognitum equidem habeo, multas alioquin impensas a te mei causa tolerari, id tamen mihi sumpsi, ut novam insuper gratiam etiam atque etiam a te contenderem. Auctor nobis fit magister, ut in coemendos eos libros, quos in schola nobis miro evolvit modo, pecuniæ aliquid collocaremus; fidemque suam obligat, fore ut, si eosdem domi etiam legendo contriverimus, magnos progressus faciamus. Uti hactenus nervis omnibus contendi, principem in schola locum vel primo proximum tenere, ita nihil tristitia majore me conficeret, atque si, ob hunc librorum defectum, priorum laudum jacturam facerem, et a condiscipulis despectui haberer. Spe autem ducor in causa tam æqua, fructus liberalitatis tuæ uberrimos, ut pervenere, ad me perventuros.


\end{minipage}

\begin{multicols}{2}
\textit{Après avoir souligné, dans mon devoir, les expressions qui me déplaisent, je cherche à les remplacer au moyen du dictionnaire de Phraséologie. Prenons, pour commencer, le verbe, scio ; à cet article nous trouverons: compertum habeo, cognitum habeo, etc. Donc compertum vel cognitum equidem habeo multas alioquin impensas a te mei causa tolerari, en substituant ce verbe passif, forme essentiellement latine, au verbe facere, gallicisme français, dans le sens de facere sumptus, faire des dépenses. Nous remplacerons, toujours au moyen de la Phraséologie, ausus tamen sum par id tamen mihi sumpsi, et à la place du gallicisme petere gratiam, nous mettrons contendere gratiam. Avec le verbe sumpsi suivi du subjonctif, nous terminerons ainsi la période: Id tamen mihi sumpsi ut novam insuper gratiam etiam atque etiam a te contenderem. C’est ainsi que nous trouverons dans le Dictionnaire tous les autres changements à faire. Tantôt ce sera le substantif, tantôt l’adjectif, le verbe ou même les dérivés, qui seront l’objet des remarques de l’élève studieux, et il est évident que l’habitude prise de corriger ainsi, pendant une heure, un discours de deux pages, ne pourra manquer de produire, en peu de temps, des fruits merveilleux. L’enfant sera heureux de reconnaître dans son latin quelques unes de ces tournures qu’on lui a appris à admirer, et il fera des progrès rapides, s’il parvient à se garder contre l’écueil opposé, de vouloir tout farcir d’élégances plus ou moins bien employées “ ne omnia Elegantiis farcire velit, sed ut condimentum parce inspergat et cum judicio, quo discernat quæ Phrasis, cui voci ejus sententiæ aptior sit.”}

\textit{Mais, me dira-t-on peut-être, à quoi bon cet amas et d’élégances, de tours propres à la langue latine, qui ne peut servir qu’à tromper professeurs et élèves, à égarer même ces derniers dans une forêt d’expressions dont ils ne sont pas capables de sentir la justesse et la portée? Ce double reproche serait sérieux, s’il était fondé. Pourquoi, dirai-je à mon tour, recommander aux élèves de choisir les expressions, les tournures les plus latines? Pourquoi ces cahiers où l’on recueille, jour par jour, les phrases les plus saillantes, les tournures les plus autorisées de la langue de Cicéron? Pourquoi ces leçons apprises chaque jour? Pourquoi ces imitations quotidiennes du prince des écrivains latins? Peut-être sera-t-il difficile, dans les commencements de saisir toute la portée du livre que nous offrons de nouveau au public lettré, mais un peu d’expérience en montrera certainement l’utilité. Ce livre a pour lui la sanction des années, et notre époque, quoique moins enthousiaste de la langue latine que les siècles précédents, conserve cependant assez de grands et nobles esprits, pour qu’il soit inutile de prouver l’utilité de l’étude des belles-lettres pour la véritable formation de l’homme et du chrétien.}

\textit{Nous ne pouvons terminer ce travail sans indiquer les améliorations apportées à cette première édition française.}

\textit{Les textes ont été revus soigneusement d’après trois éditions différentes. Quand ils étaient obscurs ou écourtés, ils ont été rétablis dans leur intégrité au moyen de Forcellini, de Freund ou des originaux, et traduits en français. Pour plus de clarté, on a établi des différences de caractères, qui permettent de trouver rapidement ce qu’on recherche. Ainsi que dans tous les lexiques modernes, la quantité a été indiquée sur les mots, pour permettre aux élèves d’observer la règle si importante de l’accent tonique.}

\textit{Trois Index complètent l’ouvrage: dans le premier, on a signalé les barbarismes les plus fréquents que peuvent commettre les élèves, et le mot latin qui les remplace. Le second est aussi curieux qu’instructif. Il rend une quantité de mots techniques et propres aux arts, aux sciences et à la jurisprudence.}

\textit{Le troisième et le plus important est l’Index Gallico-Latinus. Les mots ont été placés par ordre alphabétique, en prenant le nom comme base du classement. Sous le nom sont placés les verbes, les adjectifs, les adverbes, et ainsi d’un seul coup-d’œil on voit les relations entre la langue française et la langue latine, les changements qui peuvent se produire dans la traduction du substantif, de l’adjectif et du verbe. Ce travail n’est pas un Dictionnaire français-latin; nous n’avons voulu faire qu’un Index, qui renvoyât au Corpus Phraseologiæ.}

\textit{Puissent ces divers travaux, entrepris pour la plus grande gloire de DIEU, édités avec le plus grand soin, répondre au but que s’est proposé l’auteur en les composant! Puissent-ils être utiles aux maîtres et aux élèves, augmenter dans les cœurs des uns et des autres l’amour de la langue latine, et avec cet amour, le désir de défendre et de protéger toujours la gardienne et la protectrice du latin, Celle qui a tant fait pour le succès des bonnes études, l’Église de DIEU!}

\textit{ Amiens, École libre de la Providence, 23 mai 1878.}

\begin{center}\textit{En la fête du bienheureux martyr André Bobola, de la Compagnie de Jésus.}\end{center}

\begin{flushright}\textit{Aug. BORGNET, S. J.}\end{flushright}

\end{multicols}
\end{otherlanguage}